\section{A Sensoria okoszokni integrációja}

Az alkalmazás egyik központi eleme a Sensoria okoszokniból érkező szenzoradatok kezelése. A
zokni külső mérőeszközként kapcsolódik az alkalmazáshoz, és a gyakorlatok végrehajtása közben
folyamatosan küldi az adatokat. A feladat ezért nemcsak az adatok fogadása volt, hanem azok
értelmezése és olyan formába alakítása is, amely az alkalmazás további részei számára már
feldolgozható.

A kommunikáció Bluetooth Low Energy-kapcsolaton keresztül történik. A kapcsolat kezeléséért egy külön
komponens felel, amely az okoszoknihoz való csatlakozást, a kapcsolat állapotának figyelését és az
adatcsomagok fogadását végzi. Ez azért különösen fontos, mert a Bluetooth-kommunikáció platformfüggő
és állapotérzékeny folyamat: kezelni kell a csatlakozás létrejöttét, az esetleges kapcsolatvesztést,
valamint azt is, hogy az adatok csak sikeres kapcsolat esetén kerüljenek feldolgozásra.

Az alkalmazás egy \code{ScanFilter} segítségével a Sensoria szolgáltatáshoz tartozó UUID alapján szűri
az elérhető eszközöket. Emellett az eszköz neve alapján is ellenőrzi, hogy valóban a zokniról van-e szó.
Amikor megtalálja a megfelelő eszközt, leállítja a keresést, majd \code{connectGatt} hívással GATT-
kapcsolatot hoz létre a zoknival. A GATT, vagyis a Generic Attribute Profile, a Bluetooth Low Energy
eszközök közötti adatkommunikáció egyik alapvető modellje. Ebben a felépítésben az okoszokni
szerverként működik, amely különböző szolgáltatásokat és azokon belül characteristic-eket tesz
elérhetővé. Az alkalmazás kliensként kapcsolódik ehhez a szerverhez, majd a megfelelő szolgáltatáson
keresztül fogadja a mérési adatokat.

A kapcsolat felépítése után indul a BLE service discovery folyamat. Ennek során az alkalmazás
megkeresi a Sensoria szolgáltatáshoz tartozó service UUID-t, majd azon belül kiválasztja azt a
characteristic-et, amely támogatja a \code{NOTIFY} vagy \code{INDICATE} tulajdonságokat. A
characteristic-ek a BLE szolgáltatáson belüli adattároló vagy adatküldő pontokként értelmezhetők.
A \code{NOTIFY} azt jelenti, hogy az eszköz új adat érkezésekor automatikusan értesítést küld a
kliensnek, visszaigazolás nélkül. Az \code{INDICATE} ehhez hasonlóan működik, de a kliens oldaláról
visszaigazolást is igényel. Erre azért van szükség, mert a szenzoradatok nem egyszeri lekérdezéssel
érkeznek, hanem a zokni folyamatosan küldi őket a mérés során.

Az értesítések fogadásához az alkalmazás engedélyezi a kiválasztott characteristic notification
funkcióját. Ehhez a Client Characteristic Configuration Descriptor, röviden CCCD, beállítása szükséges.
Ez azt határozza meg, hogy a kliens kér-e értesítéseket vagy indikációkat az adott characteristic
változásairól. A beállítás után a zokni képes folyamatosan továbbítani az új
szenzoradatokat az alkalmazás felé.

A beérkező adatok kezelése callback-alapúan történik. Amikor a zokni új adatcsomagot küld,
az Android rendszer meghívja az \code{onCharacteristicChanged} függvényt. A megvalósítás figyelembe
veszi, hogy Android 13-tól megváltozott ennek a függvénynek a szignatúrája, ezért a kódban a régebbi
és az újabb API-verzióhoz tartozó változat is szerepel. Mindkét esetben a beérkező bájttömb kerül
továbbításra a csomagfeldolgozó logikának.

A szenzoradatok továbbítása az alkalmazás többi része felé Kotlin Flow segítségével valósul meg.
A \code{connectAndListen} függvény egy \code{Flow<SensoriaData>} típust ad vissza, amelybe a beérkező
és sikeresen dekódolt mérési adatok kerülnek. A callback-alapú Android BLE működés és a
coroutine-alapú alkalmazáslogika összekapcsolására a megvalósítás \code{callbackFlow}-t használ.
Ez lehetővé teszi, hogy a Bluetooth callbackekből érkező események reaktív adatfolyamként
jelenjenek meg az alkalmazás többi komponense számára. A kapcsolat lezárásakor az \code{awaitClose}
blokk gondoskodik a scan leállításáról, a GATT-kapcsolat bontásáról és az erőforrások felszabadításáról.

A nyers adatcsomagok dekódolását külön logika végzi. A Sensoria okoszokniból érkező bájtok nem
közvetlenül értelmezhetőek, ezért az alkalmazás bitenként olvassa ki belőlük a szükséges mezőket.
Ehhez a \code{BitReader} segédosztályt használja, amely egy bájttömbből megadott számú bitet tud
sorban kiolvasni. A dekódolás során az alkalmazás először átugorja a csomag fejlécéhez és metaadataihoz
tartozó mezőket, például az üzenettípust, a kontextust és az időzítési mezőket, majd kiolvassa a mérés
szempontjából fontos értékeket.

A feldolgozás során három nyomáscsatorna és három gyorsulási komponens kerül kinyerésre.
A \code{ch0}, \code{ch1} és \code{ch2} mezők a talp különböző pontjain mért nyomásértékeket
reprezentálják, míg az \code{accelX}, \code{accelY} és \code{accelZ} a láb három tengely menti
gyorsulását tartalmazzák. A csomagban található magnetométer-adatokat az alkalmazás jelenlegi
verziója kiolvassa, de nem használja fel az eredményképzéshez. A dekódolás eredménye egy
\code{SensoriaData} objektum, amely már strukturált formában tartalmazza a további feldolgozáshoz
szükséges adatokat.

A \code{SensoriaData} adatmodell a szenzorból származó mérési értékek mellett időbélyeget is tartalmaz.
Ez azért fontos, mert a szenzoradatok sorrendje és időbeli elhelyezkedése a későbbi feldolgozás során
is lényeges. Az adatmodellben szerepel a feladathoz tartozó információ is, így az egyes mérési pontok
összekapcsolhatók azzal a gyakorlattal, amelynek végrehajtása közben keletkeztek.

A szenzoradatok az aktuális munkamenethez kapcsolódva kerülnek eltárolásra.
A mérés indításakor az alkalmazás elkezdi gyűjteni a \code{Flow}-ból érkező \code{SensoriaData}
objektumokat, majd ezeket a feladatokhoz és a munkamenethez rendeli.
A mérés lezárásakor az összegyűjtött adatok a többi munkamenet-információval együtt mentésre kerülnek.
Így egy korábbi mérés visszatöltésekor nemcsak a munkamenet alapadatai, hanem a hozzá tartozó
szenzoros mérési adatok is elérhetők.

A megvalósítás során az Android Bluetooth-jogosultságok kezelésére is figyelni kellett. Újabb Android-
verziókon a Bluetooth-eszközök kereséséhez és a kapcsolódáshoz külön engedélyek szükségesek, például
a \code{BLUETOOTH\_SCAN} és a \code{BLUETOOTH\_CONNECT}. Régebbi verziókon a BLE scan működéséhez
helymeghatározási jogosultságra is szükség lehet. Az alkalmazás ezért a Bluetooth-műveletek indítása
előtt ellenőrzi, hogy a szükséges rendszerengedélyek rendelkezésre állnak-e.
